%! Modeling with Quadratics — Unit 3, Lesson 3.7 — Warm-Up
\documentclass[10pt]{article}
\usepackage{algebra2-article}
\usepackage{algebra2-boxes}
\newcommand{\TallMath}[1]{$\displaystyle #1\rule[-1.4em]{0pt}{3.2em}$}

\begin{document}

\pageheader{Unit 3, Lesson 3.7}{Warm-Up}
\namedateperiod

\vspace{0.10in}

\begin{notesbox}{Getting Ready: Skills We'll Use Today}
\small Today you'll \emph{build} a quadratic from a story, then read the answer off its features. These three skills
are the tools --- a vertex, a zero, and a value read from the function.

\vspace{4pt}
\textbf{1. Find the vertex} (Lesson 3.1). For $f(x)=-x^2+4x+1$, use $x=-\dfrac{b}{2a}$, then evaluate.
\[
  x=\blank{1.8cm} \qquad f(x)=\blank{1.6cm} \qquad \text{vertex } \blank{2.2cm}
\]
Because the parabola opens \underline{\hspace{2.2cm}}, this vertex value is a \textbf{max\,/\,min} \emph{(circle one)}.

\vspace{4pt}
\textbf{2. Solve by factoring} (Lesson 3.2). Set each factor to zero.
\[
  -16t^2+48t=0:\ \ t=\blank{3.0cm}
\]

\vspace{4pt}
\textbf{3. Read a height model.} A ball's height (feet) after $t$ seconds is $h(t)=-16t^2+48t$.
\[
  h(0)=\blank{1.6cm}\ \text{(what it means: \blank{4.0cm})}
\]
\emph{Today: the vertex's $y$-value is the \textbf{greatest height} the ball ever reaches.}
\end{notesbox}

\end{document}
